\section{Results}
\label{sec:results}


Unfolded-point error bars give the total uncertainty from \S\ref{sec:systematics},
combining systematic universes, statistical bootstraps, and ML training-seed
variation. Published points carry the published total covariance.
Figs.~\ref{fig:ptslices}--\ref{fig:pzslices} carry statistical errors only,
as their captions specify.

\subsection{Double-differential cross section}
\label{sec:results-xsec}
Figure~\ref{fig:xsec} shows the unfolded $d^{2}\sigma/(d\pT\,d\pPar)$ and its
1D projections. The total cross section is
\SI{\sigTwoD}{cm^2/nucleon} (published: \SI{\sigTwoDpaper}{cm^2/nucleon}).

The shapes are expected for CC-inclusive scattering at
$\langle E_\nu\rangle\sim\SI{6}{GeV}$. The $\pT$ projection vanishes at
$\pT=0$ and peaks broadly at $\pT\approx0.55$--$\SI{0.70}{GeV/c}$, reflecting
the $Q^{2}\lesssim\SI{1}{GeV^2}$ scale of few-GeV scattering through
$\pT^{2}\simeq Q^{2}\,E_\mu/E_\nu$, with $E_\mu/E_\nu\sim0.6$ for these forward
muons. It falls by more than two orders of magnitude toward the
\SI{4.5}{GeV/c} edge. The $\pPar$ projection rises from the
\SI{1.5}{GeV/c} threshold to a broad $\pPar\approx3$--$\SI{4.5}{GeV/c}$ maximum
(the NuMI medium-energy flux peak at $E_\nu\approx\SI{6}{GeV}$), then falls by two
orders of magnitude by $\pPar\approx\SI{20}{GeV/c}$. The
$\theta_\mu<20^{\circ}$ acceptance empties the 2D map's high-$\pT$, low-$\pPar$
corner (19 unreported diagonal cells; \S\ref{sec:method-binning}).



\begin{figure}[H]
  \centering
  \includegraphics[width=0.62\textwidth]{MEFHC_5iter_fig13}\\
  \includegraphics[width=0.49\textwidth]{MEFHC_5iter_xsec_proj_pt}
  \includegraphics[width=0.49\textwidth]{MEFHC_5iter_xsec_proj_pz}
  \caption{Unfolded 2D cross section (top, Fig.~13 style) and its $\pT$/$\pPar$
  projections (bottom). The $\pT$ peak at ${\sim}\SI{0.6}{GeV/c}$ reflects
  the $Q^{2}\lesssim\SI{1}{GeV^2}$ momentum-transfer scale; the $\pPar$
  peak reflects the medium-energy flux maximum.}
  \label{fig:xsec}
\end{figure}


\subsection{Comparison to the published measurement}
\label{sec:results-compare}
The published and unbinned results share data and flux, detector, and generator
systematics, with no available cross-covariance. Every covariance distance and
standardized coordinate here is therefore an indicative descriptive scale;
none is a calibrated goodness-of-fit.

Over the 205 reported bins, the total cross section reproduces the published
result~\cite{MINERvA:2021owq} to \SI{1.1}{\percent} (ratio $\ratioTot$;
Table~\ref{tab:compare}). The median per-bin ratio is $\medianBinRatio$;
$\binsFive\%$ of bins agree within \SI{5}{\percent} and $\binsTen\%$ within
\SI{10}{\percent}. Indicative covariance distances per reported bin use the
published covariance alone ($\chiPaper$) or the combined covariance, including
our uncertainty ($\chiCombined$); see \S\ref{sec:results-tension} for the former.

% chi^2 caveat, moved here from sec_systematics.tex (formerly \paragraph{$\chi^{2}$ caveat.}).
The combined paper$+$ours $\chi^{2}/\mathrm{ndf}=\chiCombined$ double-counts
shared flux/GENIE/MINOS systematics. The ours-only inverse-covariance $\chi^{2}$
uses only the well-constrained subspace under a truncated-spectral convention.
We therefore quote per-bin pulls, the regularization
scan, and the truncated result together rather than treating an unconstrained
pseudo-inverse as a standalone goodness-of-fit (App.~\ref{sec:rank}).

Per-bin differences standardized by the published uncertainty
(Fig.~\ref{fig:pull}) are centered (mean $\pullMean$), with RMS $\pullRMSr$
well inside the published uncertainties.
Figures~\ref{fig:ptslices}--\ref{fig:pzslices} compare our unfold with
\emph{simulated truth} in the published panel format. Published data do not
enter this data-versus-prior comparison.

\begin{table}[h]
  \centering
  \begin{tabular}{@{}lc@{}}
    \toprule
    Quantity (205 reported bins) & Value \\
    \midrule
    $\sigma_{\mathrm{tot}}(\mathrm{ours})/\sigma_{\mathrm{tot}}(\mathrm{paper})$ & $\ratioTot$ \\
    Median bin ratio (ours/paper) & $\medianBinRatio$ \\
    Bins within \SI{5}{\percent}/\SI{10}{\percent}/\SI{20}{\percent} & $\binsFive/\binsTen/\binsTwenty\,\%$ \\
    Indicative covariance distance/bin, paper covariance & $\chiPaper$ \\
    Indicative covariance distance/bin, paper $+$ ours & $\chiCombined$ \\
    Published-$\sigma$ standardized-difference mean / RMS & $\pullMean/\pullRMS$ \\
    \bottomrule
  \end{tabular}
  \caption{Agreement with arXiv:2106.16210. Per-bin differences standardized
  by published uncertainties (Fig.~\ref{fig:pull}) are sub-$1\sigma$;
  covariance distances are descriptive scales.}
  \label{tab:compare}
\end{table}

\begin{figure}[H]
  \centering
  \includegraphics[width=0.7\textwidth]{MEFHC_5iter_pull_full}
  \caption{Per-bin standardized difference
  $(\mathrm{ours}-\mathrm{paper})/\sigma_{\mathrm{paper}}$ on the
  $14\times16$ grid and its distribution. It is a descriptive scale, not an
  independent-result pull.}
  \label{fig:pull}
\end{figure}

\begin{figure}[H]
  \centering
  \includegraphics[width=0.99\textwidth]{MEFHC_5iter_xsec_paper_pt_slices}
  \caption{Unfolded cross section (points) and GENIE 2.12.6 $+$ MINERvA Tune~v1
  \emph{simulated truth} (histogram), sliced in $\pT$ versus $\pPar$.
  The histogram uses the analysis MC truth denominator with the unfolded
  result's flux, nucleon-count and efficiency normalization%
  % (\texttt{2d-unfolding/plot\_2d\_paper\_comparison.py} builds it from
  % \texttt{hOFTruthDenom2D}; listed in App.~\ref{app:codebase})
  . The curves differ only by data versus MC prior; no published measurement is read.
  Only the \emph{panel format} follows Fig.~13 of
  Ref.~\cite{MINERvA:2021owq}. Points carry statistical errors only;
  these alone enter the annotated per-panel $\chi^{2}$, a descriptive
  scale. Comparisons \emph{to the published data} appear in Fig.~\ref{fig:pull}
  (per-bin standardized differences), Table~\ref{tab:compare}, and
  Fig.~\ref{fig:model} (projections).}
  \label{fig:ptslices}
\end{figure}

\begin{figure}[H]
  \centering
  \includegraphics[width=0.99\textwidth]{MEFHC_5iter_xsec_paper_pz_slices}
  \caption{Unfolded versus simulated truth (Fig.~\ref{fig:ptslices}), sliced
  in $\pPar$ versus $\pT$. Statistical errors only; no published measurement
  is read.}
  \label{fig:pzslices}
\end{figure}

\subsection{Uncertainty budget}
\label{sec:results-uq}
Table~\ref{tab:uq} gives the budget over the 205 reported bins, constructed as
described in \S\ref{sec:systematics}. Our \emph{standalone} combined budget matches
the published total in size --- median per-bin \SI{\uqMedian}{\percent} against the
published \SI{\uqPaper}{\percent} --- and is flux-dominated as in the publication
(Fig.~\ref{fig:uqbands}). The
systematic block dominates; statistical and ML blocks together contribute
under a tenth of the combined $\sqrt{\mathrm{Tr}\,C}$.

\begin{table}[H]
  \centering
  \begin{tabular}{@{}lccc@{}}
    \toprule
    Component & $N$ & $\sqrt{\mathrm{Tr}\,C}$ & per-bin median rel.\ $\sigma$ \\
    \midrule
    ML noise (lgbm training-seed var.) & 10 & \num{5.06e-41} & \SI{0.166}{\percent} \\
    Statistical (Poisson bootstrap) & 300 & \num{1.93e-40} & \SI{0.674}{\percent} \\
    Systematic (universe sweep $+$ \SI{1.4}{\percent} norm) & $187{+}1$ & \num{3.21e-39} & \SI{6.830}{\percent} \\
    \textbf{Combined (block sum)} & --- & \num{3.22e-39} & \SI{6.871}{\percent} \\
    Paper TotalCov (reference) & --- & \num{2.68e-39} & \SI{\uqPaper}{\percent} \\
    \bottomrule
  \end{tabular}
  \caption{2D uncertainty budget over 205 reported bins. Each component has
  a $205\times205$ covariance matrix $C$; $\sqrt{\mathrm{Tr}\,C}$ gives its
  absolute per-bin uncertainties in quadrature. The standalone combined
  median relative uncertainty (\SI{6.871}{\percent}) matches the published
  total (\SI{\uqPaper}{\percent}).}
  \label{tab:uq}
\end{table}
\subsection{Anatomy of the $\chi^{2}=\chiPaper$ tension}
\label{sec:results-tension}
Using the published uncertainties as a descriptive scale, the central values
have pull RMS $\pullRMS$ (a diagonal-only $\chi^{2}/\mathrm{ndf}$ would be
$0.37$), while the full paper-covariance distance is $\chiPaper$ per bin:
a $10\times$ inflation living entirely in the off-diagonal structure of the
published covariance. The eigenmode decomposition (Fig.~\ref{fig:tension}) shows
the indicative distance is carried by moderate-$\lambda$ modes whose
paper-covariance-standardized components have magnitude $5$--$6$
(the 10 smallest-$\lambda$ modes carry only \SI{3}{\percent}; keeping the 73
best-measured directions still gives $1.42$), so it is \emph{not} a small-$\lambda$
inversion artifact. It is \emph{shape}, not normalization
($\chi^{2}_{\mathrm{norm}}=0.10$ vs $\chi^{2}_{\mathrm{shape}}=750$), localized to
the low-$\pT$ cross-section peak ridge where the flux and muon-energy bands
dominate. The GBDT estimator gives a $\sim1$-unit methodological band
(exact-GBT $\chiPaper$ vs HistGBT $2.70\pm0.04$ vs LightGBM $2.65$).
Iteration-count effects are negligible: $5\to10$ iterations move it by $+0.01$.
See App.~\ref{sec:tension} for the full treatment and
Figs.~\ref{fig:tension}--\ref{fig:tensionmodes} for eigenmode-localization maps.
% NOTE: the estimator-band value "exact-GBT 3.66" below is the same number as
% the paper-covariance chi2 (\chiPaper); they coincide by construction.

\begin{figure}[H]
  \centering
  \includegraphics[width=0.49\textwidth]{tension_spectrum}
  \includegraphics[width=0.49\textwidth]{tension_chi2_map}
  \caption{Tension anatomy: eigenvalue spectrum with cumulative $\chi^{2}$ (left)
  and the per-bin $\chi^{2}$-contribution map (right).}
  \label{fig:tension}
\end{figure}

\begin{figure}[H]
  \centering
  \includegraphics[width=0.9\textwidth]{tension_mode_maps}
  \caption{Leading tension-carrying eigenmodes of the published covariance
  in bin space. Magnitude-$5$--$6$ standardized components occupy the low-$\pT$
  cross-section-peak ridge where flux and muon-energy bands dominate.}
  \label{fig:tensionmodes}
\end{figure}

\subsection{Comparison to GENIE MINERvA Tune v1}
\label{sec:results-model}
GENIE 2.12.6~\cite{Andreopoulos:2009rq} MINERvA Tune v1 under-predicts the total
by $\sim\SI{9}{\percent}$ ($\sigma_{\mathrm{tune}}/\sigma_{\mathrm{data}}=0.91$),
concentrated in the low-$\pT$ peak. The ancillary prediction is compared with
both the published data and our result in Fig.~\ref{fig:model}. The data
covariance strongly disfavors the tune ($\chi^{2}/\mathrm{ndf}$: data vs tune
$33.0$, ours vs tune $26.5$). Our unfold tracks the data far more closely than
the model does (Fig.~\ref{fig:modelmaps}).
The published analysis~\cite{MINERvA:2021owq} also compares with
NuWro~\cite{Golan:2012wx}, GiBUU~\cite{Buss:2011mx} and NEUT~\cite{Hayato:2009zz}.
OmniFold's per-event reweighted truth ensemble can be compared with any
such prediction in \emph{any} projection of classifier features without
re-unfolding. Adding an observable outside that space is a new
unfold, not a projection: the 3D extension (\S\ref{sec:3d}) adds $\Eavail$ as a
classifier feature and is re-unfolded accordingly.

\begin{figure}[H]
  \centering
  \includegraphics[width=0.85\textwidth]{model_comp_projections}
  \caption{Our result (dashed) on top of the published data (black); MINERvA
  Tune v1 (orange) dips below in the peak.}
  \label{fig:model}
\end{figure}

\begin{figure}[H]
  \centering
  \includegraphics[width=0.9\textwidth]{model_comp_ratio_maps}
  \caption{Bin-by-bin ratios to MINERvA Tune~v1 on the $(\pT,\pPar)$ grid:
  ours/Tune~v1 (left) and published data/Tune~v1 (right), with common colour
  scale $0.7$--$1.3$
  % (\texttt{2d-unfolding/compare\_to\_models.py}; listed in App.~\ref{app:codebase})
  (only this caption labels the panels; they have no printed titles).
  Neither map uses published data as its denominator or is flat near unity.
  Both show the same tune deficit, above unity over most of the plane and
  dipping below at the lowest $\pT$. The close maps reflect that our result departs from
  the tune in the same places and by similar amounts as the published data. Direct comparison with
  published data uses Table~\ref{tab:compare}'s median bin ratio and
  Fig.~\ref{fig:pull}'s residual map.}
  \label{fig:modelmaps}
\end{figure}
